\begin{center}
\captionof{table}{$L{=}256$ \rs{} composition (125M, 3-seed mean PPL). The
tested $\tau{=}4.0$ setting has the lowest mean at every reported ratio.}
\label{tab:pe-yarn-l256}
\small
\setlength{\tabcolsep}{6pt}
\begin{tabular}{@{}lcccc@{}}
\toprule
Method (\(L{=}256\)) & 4$\times$ & 8$\times$ & 16$\times$ & 32$\times$ \\
\midrule
Geo+\rs{} & 60.7 & 97.5 & 190.4 & 260.2 \\
\evq{} ($\tau{=}2$)+\rs{} & 51.6 & 66.7 & 107.9 & 141.7 \\
\evq{} ($\tau{=}4$)+\rs{} & \textbf{50.9} & \textbf{60.9} & \textbf{84.1} & \textbf{99.6} \\
\bottomrule
\end{tabular}
\end{center}

\begin{center}
  \includegraphics[width=0.96\linewidth]{fig_range_composition_125m.pdf}
  \captionof{figure}{\textbf{The composition gap widens with extrapolation.}
  At 125M, the tested $\tau{=}4$ allocation remains lowest from $4\times$ to
  $32\times$ under the same \rs{} rule.}
  \label{fig:range-composition-125m}
\end{center}

\paragraph{Reading the composition.}
The tested $\tau{=}4$ composition lowers PPL relative to Geo+\rs{} by
$16\%/38\%/56\%/62\%$ at $4\times/8\times/16\times/32\times$. In this
protocol, the same index-space rule produces different composition effects on
the two learned spectral substrates.
